% Part: many-valued-logic
% Chapter: syntax-and-semantics
% Section: matrices

\documentclass[../../../include/open-logic-section]{subfiles}

\begin{document}

\olfileid{mvl}{syn}{mat}

\olsection{Matriks}

Logika mawa akeh nilai ditemtokake dening basane, himpunan nilai
kabenerane~$V$, subhimpunan nilai kabeneran pinilih, lan fungsi
kabeneran kanggo saben konektife. Bebarengan, unsur-unsur iki
diarani \emph{matriks}.

\begin{defn}[Matriks]
\ollabel{defn:matrix}
\emph{Matriks} kanggo logika~$\Log L$ dumadi saka:
\begin{enumerate}
\item himpunan konektif kang mbentuk basa~$\Lang L$;
\item himpunan nilai kabeneran $V \neq \emptyset$;
\item himpunan nilai kabeneran pinilih $V^+ \subseteq V$;
\item kanggo saben konektif kanthi $n$ panggonan, yaiku $\star$, ing
$\Lang L$, fungsi kabeneran~$\tf{\star} : V^n \to V$.
Yen $n = 0$, $\tf{\star}$ mung sawijining unsur saka~$V$.
\end{enumerate}
\end{defn}

\begin{ex}
Matriks kanggo logika klasik~\LogCL{} dumadi saka:
\begin{enumerate}
  \item Basa proposisional baku $\Lang L_0$ kang ngemot
  $\lfalse$, $\lnot$, $\land$, $\lor$, $\lif$.
  \item Himpunan nilai kabeneran $V = \{\True, \False\}$.
  \item $\True$ mung siji-sijine nilai pinilih, yaiku
  $V^+ = \{\True\}$.
  \item Kanggo $\lfalse$, kita duwe $\tf{\lfalse} = \False$.
  Fungsi kabeneran liyane diwenehake dening tabel kabeneran
  kang lumrah (delengen \olref{fig:tf-CL}).
\end{enumerate}
\begin{figure}
  \begin{center}
      \begin{tabular}{c|c} 
        $\tf{\lnot}$ & \\ 
        \hline  
        $\True$ & $\False$ \\ 
        $\False$ & $\True$ 
      \end{tabular}
      \quad
      \begin{tabular}{c|cc} 
        $\tf{\land}$ & $\True$ & $\False$ \\ 
        \hline 
        $\True$ & $\True$ & $\False$ \\ 
        $\False$ & $\False$ & $\False$ 
      \end{tabular}
      \quad
      \begin{tabular}{c|cc} 
        $\tf{\lor}$ & $\True$ & $\False$ \\ 
        \hline 
        $\True$ & $\True$ & $\True$ \\ 
        $\False$ & $\True$ & $\False$ 
      \end{tabular}
      \quad
      \begin{tabular}{c|cc} 
        $\tf{\lif}$ & $\True$ & $\False$ \\ 
        \hline 
        $\True$ & $\True$ & $\False$ \\ 
        $\False$ & $\True$ & $\True$ 
      \end{tabular}
    \end{center} 
    \caption{Fungsi kabeneran kanggo logika klasik~$\LogCL$.}
    \ollabel{fig:tf-CL}
  \end{figure}
  \end{ex}

\end{document}
